\section*{\S\,5. Машинный способ задания функций. Коммутативные диаграммы}
\addcontentsline{toc}{section}{\S\,5. Машинный способ задания функций. Коммутативные диаграммы}
\markright{\S\,5. Машинный способ задания функций}
\subsection*{Условия}
1. Доказать, что множество $C^m(\R^n)$ ($\forall n\in\N$) $m$ раз везде
дифференцируемых функций (всевозможного
конечного числа аргументов) замкнуто. Верно ли это утверждение и при
$m=\infty$?

2. Доказать, что множество $C^A(\R^n)$ ($\forall n\in\N$) везде
аналитических функций (т.\,е. раскладывающихся в ряд Тейлора с ненулевым
радиусом сходимости в любой точке $\R^n$) замкнуто.

3. Показать, что замкнутым множеством является множество $\varepsilon_n$
$n$ раз непрерывно дифференцируемых элементарных функций.

4. Построить полную систему в $\varepsilon_n$ (пример
3),\notefn{по-видимому, имеется в виду упражнение~3, в котором определено
множество $\varepsilon_n$.} содержащую конечное число функций (не считая
констант из $\R$).

5. Пусть $\varphi_i\in C(\R^{n_i})$ ($i=1,\dots,s$), где $C(\R^k)$ "---
множество всех непрерывных в $\R^k$ функций. Пусть $f$ есть суперпозиция
функций $\varphi_i$, имеющая структуру
$f=\chi[\varphi_1(x_{11},\dots,x_{n_11}),\dots,\varphi_s(x_{1s},\dots,x_{n_ss})]$.
Очевидно, что $f$ как суперпозиция непрерывных функций есть непрерывная
функция вида $f(x_1,\dots,x_n)$, где $x_1,\dots,x_n$ получается из строки
аргументов $x_{11},\dots,x_{n_11},\dots,x_{1s},\dots,x_{n_ss}$ путем
вычеркивания повторно встречающихся букв. Таким образом, приходим
к отображению $\chi^*\colon\R^{n_1}\times\dots\times\R^{n_s}\to\R^n$ (где
$\chi^*(\R^{n_1}\times\dots\times\R^{n_s})\subset C(\R^n)$ "--- образ этого
отображения), вполне определяемому структурой
$\chi$.\notefn{по смыслу отображение $\chi^*$ сопоставляет набору функций
$(\varphi_1,\dots,\varphi_s)$ функцию $f$, т.\,е. действует из
$C(\R^{n_1})\times\dots\times C(\R^{n_s})$ в $C(\R^n)$; так же следует
понимать и запись $\chi(\R^{n_1}\times\dots\times\R^{n_s})$ в следующем
абзаце "--- как образ этого отображения.}

Для каких структур $\chi(\R^{n_1}\times\dots\times\R^{n_s})=C(\R^n)$?
Привести пример структуры, для которой такое равенство не достижимо.
Выяснить, в частности, всякая ли непрерывная функция $f(x,y,z,t)$ может быть
представлена в виде $f(x,y,z,t)\equiv\varphi(x,\varphi(y,\varphi(z,t)))$,
где $\varphi(x,y)\in C(\R^2)$? В упоминавшейся в \S\,4 (пример~2)
13~проблеме Гильберта речь шла о представлении непрерывных функций трех и
большего числа аргументов в виде суперпозиций непрерывных функций двух
аргументов; при этом разрешалось использовать формулы любой структуры.
Здесь же ставится вопрос так: нельзя ли сократить множество структур,
используемых для представления непрерывных функций трех и более аргументов
через непрерывные функции двух аргументов? Может быть существует такая
одна структура, которой будет достаточно для этой цели?

6. Вопросы, поставленные в задаче~5, могут быть перенесены и на другие
функциональные классы ($C^m(\R^n)$, $H^s(\R^n)$ и~др.).

\subsection*{Решения}

Напомним определения из \S\,4, которые понадобятся в упражнениях 1--4.
Множество $\M(H)$ \emph{суперпозиций} системы $H$ включает: а)~отображения
системы $H$; б)~тождественное отображение; в)~пары $(\varphi,\psi)$
отображений из $\M(H)$ (компоненты пары "--- компоненты $\varphi$ и $\psi$);
г)~композиции $\varphi\circ\psi$ отображений из $\M(H)$. Иначе говоря,
суперпозиции получаются из отображений системы $H$ и тождественного
отображения конечным числом применений правил в) и~г). Как замечено
в \S\,4, при этом допускается произвольное переименование переменных,
в том числе их отождествление: суперпозицией там названа, например, функция
$f[\varphi(x_1,x_2),x_1,f(x_1,x_1,x_1)]$. Множество $\M_0$ называется
\emph{замкнутым}, если $\M(H)\subset\M_0$ для всякой системы $H\subset\M_0$
(знак $\subset$ в книге допускает равенство). Правило в) дает векторные
отображения; для множества числовых функций включение $\M(H)\subset\M_0$
будем понимать покомпонентно: каждая компонента каждого отображения из
$\M(H)$ лежит в $\M_0$.

\begin{solution}[1 (\S\,5)]
\textbf{Что означает $C^m$.} В \S\,4 множество $C^k(\R^n)$ $k$ раз
непрерывно дифференцируемых везде функций названо среди примеров замкнутых
множеств; упражнение просит это доказать. Будем пользоваться таким
определением по индукции: $C^0(\R^n)$ "--- непрерывные на $\R^n$ функции;
функция принадлежит $C^{m+1}(\R^n)$, если у нее во всех точках существуют
все частные производные первого порядка и все они принадлежат $C^m(\R^n)$.
Иначе говоря, $C^m$ "--- функции, у которых существуют и непрерывны все
частные производные до порядка $m$ включительно. Отсюда $C^{m+1}\subset C^m$:
функция с непрерывными частными производными первого порядка
дифференцируема (теорема из курса анализа) и потому непрерывна; применяя
это к производным порядка не выше $m$ функции класса $C^{m+1}$, видим, что
они непрерывны. Наконец, $C^\infty=\bigcap_mC^m$. О буквальном прочтении
слов <<$m$ раз везде дифференцируемых>> (без слова <<непрерывно>>) сказано
в конце решения.

\textbf{Проба.} Посмотрим на примере, как устроена гладкость суперпозиции.
Функция $F(u)=u\,|u|$ имеет непрерывную производную $F'(u)=2|u|$, а второй
производной в нуле у нее нет: $F\in C^1(\R)$, $F\notin C^2(\R)$. Подставим
$u=\varphi(x,y)=x+y^2$. По правилу дифференцирования сложной функции у
$g(x,y)=F(\varphi(x,y))=(x+y^2)\,|x+y^2|$
\[
\frac{\partial g}{\partial x}=F'(\varphi)\,\frac{\partial\varphi}{\partial x}=2|x+y^2|,\qquad
\frac{\partial g}{\partial y}=F'(\varphi)\,\frac{\partial\varphi}{\partial y}=4y\,|x+y^2|;
\]
обе производные непрерывны, и $g\in C^1(\R^2)$. На параболе $x=-y^2$ у
производной $\partial g/\partial x=2|x+y^2|$ нет производной по $x$, так что
$g\notin C^2$: суперпозиция ровно так же гладка, как <<худшая>> из
составляющих функций. Важна форма вычисления: производная суперпозиции
снова составлена из суперпозиции ($F'\circ\varphi$), но уже функций на одну
ступень менее гладких, и из производных $\varphi$. Это подсказывает индукцию
по $m$.

\textbf{Лемма о суперпозиции.} Докажем индукцией по $m$ утверждение
\begin{keystmt}
$P(m)$: при любых $n,p\ge1$, любой $F\in C^m(\R^p)$ и любых
$\varphi_1,\dots,\varphi_p\in C^m(\R^n)$ функция
$g(x)=F(\varphi_1(x),\dots,\varphi_p(x))$ принадлежит $C^m(\R^n)$.
\end{keystmt}
Утверждение сформулировано сразу для всех $n$ и $p$: в шаге индукции оно
понадобится при других значениях $p$.

\emph{База}, $m=0$: суперпозиция непрерывных функций непрерывна. Если
$x^{(k)}\to x$, то $\varphi_i(x^{(k)})\to\varphi_i(x)$ при каждом $i$, т.\,е.
точки $\varphi(x^{(k)})$ сходятся к $\varphi(x)$ в $\R^p$, и по
непрерывности $F$ имеем $g(x^{(k)})\to g(x)$.

\emph{Шаг.} Пусть $P(m)$ верно; докажем $P(m+1)$. Пусть $F\in C^{m+1}(\R^p)$,
$\varphi_i\in C^{m+1}(\R^n)$. Так как $m+1\ge1$, у функций $F$, $\varphi_i$
есть непрерывные частные производные первого порядка, а такие функции
дифференцируемы. Поэтому применимо правило дифференцирования сложной
функции:
\begin{equation*}\tag{1}
\frac{\partial g}{\partial x_j}(x)=\sum_{k=1}^{p}\frac{\partial F}{\partial y_k}\bigl(\varphi(x)\bigr)\,
\frac{\partial\varphi_k}{\partial x_j}(x),\qquad j=1,\dots,n .
\end{equation*}
Разберем правую часть. Производная $\partial F/\partial y_k$ лежит
в $C^m(\R^p)$, потому что $F\in C^{m+1}$; функции $\varphi_1,\dots,\varphi_p$
лежат в $C^{m+1}\subset C^m$; по предположению $P(m)$ функция
$(\partial F/\partial y_k)\circ\varphi$ лежит в $C^m(\R^n)$. Множители
$\partial\varphi_k/\partial x_j$ лежат в $C^m(\R^n)$, потому что
$\varphi_k\in C^{m+1}$. Наконец, вся сумма в (1) "--- это функция
$S(u_1,\dots,u_p,v_1,\dots,v_p)=\sum_ku_kv_k$ от этих $2p$ функций, а $S$
"--- многочлен, так что $S\in C^\infty(\R^{2p})\subset C^m(\R^{2p})$; по $P(m)$
(с $2p$ вместо $p$) правая часть (1) лежит в $C^m(\R^n)$. Итак, у $g$ всюду
есть все первые частные производные, и все они из $C^m(\R^n)$, т.\,е.
$g\in C^{m+1}(\R^n)$. Утверждение $P(m+1)$ доказано.

\textbf{От леммы к замкнутости.} Пусть $H\subset C^m$ "--- произвольная
система (функции из $H$ могут зависеть от разного числа переменных).
Назовем \emph{рангом} отображения из $\M(H)$ наименьшее число применений
правил в), г) и переименований переменных, с помощью которого оно
получается из функций системы $H$ и тождественного отображения. Докажем
индукцией по рангу $r$ утверждение
\begin{keystmt}
$Q(r)$: у отображения из $\M(H)$ ранга $r$ все компоненты лежат в $C^m$.
\end{keystmt}
\emph{База}, $r=0$: это функции из $H$ (они лежат в $C^m$ по условию)
и тождественное отображение $y=x$ (многочлен).

\emph{Шаг.} Пусть $Q(0),\dots,Q(r)$ верны и отображение имеет ранг $r+1$.
Последняя операция применена к одному или двум отображениям меньшего ранга,
причем их ранги могут быть любыми, не большими $r$, "--- поэтому нужна
полная индукция. Возможны три случая.
\begin{itemize}[leftmargin=1.2em,itemsep=2pt]
\item Пара $(\varphi,\psi)$: ее компоненты "--- компоненты $\varphi$ и $\psi$,
  и они лежат в $C^m$ по предположению индукции.
\item Композиция $\varphi\circ\psi$: каждая ее компонента имеет вид
  $F(\psi_1,\dots,\psi_p)$, где $F$ "--- компонента $\varphi$, а $\psi_k$ "---
  компоненты $\psi$; все они лежат в $C^m$, и по лемме
  $F(\psi_1,\dots,\psi_p)\in C^m$.
\item Переименование переменных: компонента получается из компоненты $F$
  прежнего отображения подстановкой переменных вместо аргументов, например
  $F(x_1,x_1,x_3)$ или $F(x_2,x_1)$. Это суперпозиция $F$ и функций вида
  $x_i$ (многочленов), и по лемме она лежит в $C^m$.
\end{itemize}
Итак, $\M(H)\subset C^m$ для всякой системы $H\subset C^m$, т.\,е. множество
$C^m$ замкнуто.

\textbf{Случай $m=\infty$.} Пусть $H\subset C^\infty$. Тогда $H\subset C^m$
при каждом $m$, и по доказанному все компоненты всех отображений из $\M(H)$
лежат в $C^m$ при каждом $m$, т.\,е. в $C^\infty$. Это частный случай
теоремы~1 \S\,4: пересечение замкнутых множеств $C^0,C^1,C^2,\dots$ замкнуто
(доказательство теоремы~1 годится для любого числа множеств).

\textbf{О формулировке.} В условии сказано <<$m$ раз везде
дифференцируемых>>, без слова <<непрерывно>>. Если понимать класс
буквально "--- как функции, имеющие дифференциалы до порядка $m$, но,
возможно, с разрывными старшими производными, "--- доказательство то же
самое. Обозначим этот класс $D^m$: $D^0$ "--- все функции, $D^{m+1}$ "---
дифференцируемые функции, все первые частные производные которых лежат
в $D^m$. База индукции тогда верна без доказательства (всякая функция лежит
в $D^0$), включение $D^{m+1}\subset D^m$ проверяется той же индукцией, а
в шаге формула (1) применима, потому что функции из $D^{m+1}$
дифференцируемы по определению. Так что и в этом прочтении множество
замкнуто.

\answer Множество $C^m$ замкнуто при каждом $m=0,1,2,\dots$: по формуле
(1) производная суперпозиции составлена из суперпозиций функций на одну
ступень менее гладких, и индукция по $m$ показывает, что суперпозиция
функций класса $C^m$ снова принадлежит $C^m$. Утверждение верно и при
$m=\infty$: $C^\infty=\bigcap_mC^m$, а пересечение замкнутых множеств
замкнуто.
\end{solution}

\begin{solution}[2 (\S\,5)]
\textbf{Что означает аналитичность.} Для функций многих переменных удобны
мультииндексы: $\alpha=(\alpha_1,\dots,\alpha_n)$ с целыми
$\alpha_j\ge0$, $|\alpha|=\alpha_1+\dots+\alpha_n$,
$x^\alpha=x_1^{\alpha_1}\cdots x_n^{\alpha_n}$. Функция $f$
\emph{аналитична в точке} $x^0$, если найдутся число $r>0$ и коэффициенты
$a_\alpha$ такие, что
\begin{gather*}
\sum_\alpha|a_\alpha|\,r^{|\alpha|}<\infty,\\
f(x)=\sum_\alpha a_\alpha(x-x^0)^\alpha\quad\text{при }|x_j-x^0_j|\le r,\ j=1,\dots,n .
\end{gather*}
В кубе $|x_j-x^0_j|\le r$ члены ряда по модулю не больше $|a_\alpha|r^{|\alpha|}$,
так что ряд сходится абсолютно и его сумма не зависит от порядка
суммирования; это ряд Тейлора функции $f$ в точке $x^0$. При $n=1$ условие
равносильно тому, что ряд Тейлора имеет ненулевой радиус сходимости
и в окрестности точки сходится к $f$; при $n>1$ абсолютную сходимость мы
включаем в определение, чтобы не оговаривать порядок суммирования кратного
ряда. $C^A(\R^n)$ "--- функции, аналитические в каждой точке $\R^n$.

\textbf{Проба.} Аналитическая функция в окрестности каждой точки "--- сумма
степенного ряда. Естественно ожидать, что ряд для суперпозиции получается
подстановкой одного ряда в другой. Проделаем это для $e^{\sin x}$ в нуле:
подставим $y=\sin x=x-\frac{x^3}{6}+\dots$ в
$e^y=1+y+\frac{y^2}{2}+\frac{y^3}{6}+\frac{y^4}{24}+\dots$ и соберем
члены до $x^4$. Так как $y^2=x^2-\frac{x^4}{3}+\dots$, $y^3=x^3+\dots$,
$y^4=x^4+\dots$, получаем
\begin{align*}
e^{\sin x}&=1+x+\frac{x^2}{2}+\Bigl(-\frac16+\frac16\Bigr)x^3
+\Bigl(-\frac16+\frac1{24}\Bigr)x^4+\dots\\
&=1+x+\frac{x^2}{2}-\frac{x^4}{8}+\dots
\end{align*}
Здесь два места требуют обоснования. Во-первых, мы раскрыли скобки
в бесконечном числе слагаемых и перегруппировали члены по степеням $x$;
это законно для абсолютно сходящихся рядов. Во-вторых, подстановка имеет
смысл, только пока значения внутреннего ряда остаются там, где сходится
внешний. Например, $\frac1{1-y}=\sum_ky^k$ сходится лишь при $|y|<1$;
подстановка $y=2x$ дает ряд $\sum_k2^kx^k$, сходящийся лишь при
$|x|<\frac12$. Поэтому ряд для суперпозиции надо рассматривать в
достаточно малой окрестности, где внутренние функции мало отличаются от
своих значений в центре. (Оба вычисления проверены программой
\texttt{artifacts/checks/section5.py}.)

\textbf{Лемма о подстановке.} \emph{Если функции $\varphi_1,\dots,\varphi_p$
аналитичны в точке $x^0$, а $F$ аналитична в точке
$y^0=(\varphi_1(x^0),\dots,\varphi_p(x^0))$, то $F(\varphi_1,\dots,\varphi_p)$
аналитична в точке $x^0$.}

Сдвигом переменных сведем дело к $x^0=0$, $y^0=0$: функции
$\tilde\varphi_k(x)=\varphi_k(x^0+x)-y^0_k$ и $\tilde F(y)=F(y^0+y)$
аналитичны в нуле (их ряды "--- те же ряды, записанные по степеням $x$ и $y$),
$\tilde\varphi_k(0)=0$ и $\tilde F(\tilde\varphi(x))=F(\varphi(x^0+x))$.
Итак, пусть
\begin{gather*}
F(y)=\sum_\alpha a_\alpha y^\alpha\ \ (|y_k|\le r),\qquad A=\sum_\alpha|a_\alpha|r^{|\alpha|}<\infty;\\
\varphi_k(x)=\sum_\beta b_{k\beta}x^\beta\ \ (|x_j|\le\rho_0),\qquad B_k=\sum_\beta|b_{k\beta}|\rho_0^{|\beta|}<\infty,
\end{gather*}
причем свободные члены $b_{k0}=\varphi_k(0)=0$.

\emph{Выбор окрестности.} Положим $M_k(\rho)=\sum_\beta|b_{k\beta}|\rho^{|\beta|}$
при $0<\rho\le\rho_0$. В этой сумме $|\beta|\ge1$ (свободного члена нет),
поэтому $\rho^{|\beta|}=\rho_0^{|\beta|}(\rho/\rho_0)^{|\beta|}\le\rho_0^{|\beta|}\,\rho/\rho_0$
и $M_k(\rho)\le B_k\,\rho/\rho_0$. Возьмем $\rho\in(0,\rho_0]$ настолько
малым, чтобы $B_k\rho/\rho_0\le r$ при всех $k$; тогда $M_k(\rho)\le r$. При
$|x_j|\le\rho$ имеем $|\varphi_k(x)|\le M_k(\rho)\le r$: точка $\varphi(x)$
лежит там, где ряд для $F$ сходится, и
\[
F(\varphi(x))=\sum_\alpha a_\alpha\,\varphi_1(x)^{\alpha_1}\cdots\varphi_p(x)^{\alpha_p}.
\]

\emph{Раскрытие скобок.} Каждое произведение
$\varphi_1^{\alpha_1}\cdots\varphi_p^{\alpha_p}$ "--- произведение $|\alpha|$
абсолютно сходящихся рядов. Перемножив их почленно, получим ряд из
одночленов $c\,x^\gamma$; при $|x_j|\le\rho$ сумма модулей его членов
равна произведению сумм модулей членов сомножителей и потому не больше
$M_1(\rho)^{\alpha_1}\cdots M_p(\rho)^{\alpha_p}\le r^{|\alpha|}$.
Значит, если подставить эти разложения во внешний ряд, сумма модулей всех
получившихся членов не больше $\sum_\alpha|a_\alpha|r^{|\alpha|}=A<\infty$.
Абсолютно сходящийся кратный ряд можно суммировать в любом порядке и с любой
группировкой; собрав вместе члены с одинаковыми одночленами $x^\gamma$,
получаем
\[
F(\varphi(x))=\sum_\gamma c_\gamma x^\gamma\quad(|x_j|\le\rho),\qquad
\sum_\gamma|c_\gamma|\,\rho^{|\gamma|}\le A<\infty .
\]
Значит, $F\circ\varphi$ аналитична в нуле. Коэффициенты $c_\gamma$
получаются ровно так, как в пробе, "--- формальной подстановкой рядов.

\textbf{Замкнутость.} Дальше рассуждение то же, что в решении упражнения~1,
с леммой о подстановке вместо леммы о суперпозиции. Пусть $H\subset C^A$.
Индукцией по рангу отображения из $\M(H)$ (числу применений правил в), г)
и переименований переменных) показываем, что все его компоненты аналитичны
в каждой точке. Функции из $H$ аналитичны по условию; тождественное
отображение $y=x$ "--- многочлен, а многочлен аналитичен всюду (его ряд по
степеням $x-x^0$ конечен); компоненты пары "--- компоненты ее членов;
компонента композиции $\varphi\circ\psi$ имеет вид $F(\psi_1,\dots,\psi_p)$
и аналитична в каждой точке $x^0$ по лемме ($F$ аналитична в точке
$\psi(x^0)$, так как аналитична всюду); переименование переменных "---
подстановка многочленов $x_i$, и аналитичность сохраняется по той же лемме.

\answer Множество $C^A$ везде аналитических функций замкнуто: в окрестности
каждой точки ряд для суперпозиции получается подстановкой одних рядов
в другой, и эта подстановка законна, потому что в достаточно малом кубе
сумма модулей членов внутренних рядов (без свободных членов) не превосходит
радиуса $r$, в котором абсолютно сходится внешний ряд.
\end{solution}

\begin{solution}[3 (\S\,5)]
\textbf{Что такое $\varepsilon_n$.} Элементарные функции "---
это $H_\varepsilon$-реализуемые функции, т.\,е. суперпозиции функций системы
(1.9) (\S\,4). Многие из них определены не везде: $\ln x$ "--- лишь при $x>0$,
$x/y$ "--- при $y\ne0$. Множество $\varepsilon_n$ будем понимать как множество
элементарных функций $f(x_1,\dots,x_k)$ (при любом $k$), которые определены
во всем $\R^k$ и имеют там непрерывные частные производные до порядка $n$
включительно. Другое возможное прочтение разобрано в конце решения.

\textbf{Пробы.} Функция $\ln(1+x^2)$ определена при всех $x$ и бесконечно
дифференцируема "--- она лежит в $\varepsilon_n$ при любом $n$. Функция
$|x|^3=\bigl(\sqrt{x^2}\bigr)^3$ (о записи $|t|=\sqrt{t^2}$ см. \S\,4)
определена везде, ее производные $3x|x|$ и $6|x|$ непрерывны, а третьей
производной в нуле нет: $|x|^3\in\varepsilon_2$, $|x|^3\notin\varepsilon_3$.
Их суперпозиции, например $|\sin x|^3$ или $\bigl|\ln(1+x^2)-1\bigr|^3$,
снова элементарны (запись получается подстановкой одной записи в другую)
и снова лежат в $\varepsilon_2$. Видно, что принадлежность $\varepsilon_n$
складывается из двух свойств, и каждое из них сохраняется при суперпозиции.

\textbf{Решение.} Пусть $H\subset\varepsilon_n$ и $f$ "--- компонента
отображения из $\M(H)$.

(а) \emph{$f$ элементарна.} Функция $f$ получена из функций системы $H$
правилами образования суперпозиций. Заменим каждую функцию из $H$ ее записью
через функции системы (1.9) "--- получится запись $f$ через функции системы
(1.9), т.\,е. $f\in\M(H_\varepsilon)$. Иначе говоря, суперпозиция суперпозиций
системы $H_\varepsilon$ сама является ее суперпозицией: если
$H\subset\M(H_\varepsilon)$, то $\M(H)\subset\M(H_\varepsilon)$ (доказано
в решении упражнения~4 \S\,4, шаг~2, индукцией по числу применений правил
в), г)).

(б) \emph{$f$ определена всюду и имеет непрерывные производные до порядка
$n$.} Функции системы $H$ определены всюду и лежат в $C^n$, а множество
$C^n$ замкнуто (упражнение~1); поэтому и $f\in C^n$ (в частности, $f$
определена всюду).

Значит, $f\in\varepsilon_n$, т.\,е. $\M(H)\subset\varepsilon_n$ для всякой
системы $H\subset\varepsilon_n$. По существу это теорема~1 \S\,4:
$\varepsilon_n$ "--- пересечение двух замкнутых множеств, множества
элементарных функций и множества $C^n$.

\textbf{Другое прочтение.} Слова <<$n$ раз непрерывно дифференцируемых>>
можно понять и как <<имеющих непрерывные производные до порядка $n$ в своей
области определения>>, если эта область открыта; так понимаемое
$\varepsilon_n$ содержит, например, $\ln x$ (при $x>0$). Утверждение остается
верным и тогда. Лемма о суперпозиции из решения упражнения~1 вместе
с доказательством дословно переносится на функции, заданные на открытых
множествах: если $\varphi_1,\dots,\varphi_p$ имеют непрерывные производные до
порядка $n$ на открытом множестве $V$, а $F$ "--- на открытом множестве $U$,
то суперпозиция $F(\varphi_1,\dots,\varphi_p)$ имеет их на множестве тех
$x\in V$, для которых $(\varphi_1(x),\dots,\varphi_p(x))\in U$. Это множество
и есть область определения суперпозиции, и оно открыто, потому что прообраз
открытого множества при непрерывном отображении открыт.

\answer Множество $\varepsilon_n$ замкнуто: суперпозиция элементарных
функций элементарна (ее запись получается подстановкой записей), а
суперпозиция всюду определенных функций класса $C^n$ снова всюду определена
и принадлежит $C^n$ по упражнению~1.
\end{solution}

\begin{solution}[4 (\S\,5)]
\textbf{Что требуется.} По определению \S\,4 система $H$ \emph{полна} по
отношению к множеству $\M_0$, если $\M_0\subset\M(H)$, т.\,е. всякая функция
из $\M_0$ есть суперпозиция функций системы $H$. Сами функции системы при
этом не обязаны лежать в $\M_0$. Нужно указать полную по отношению
к $\varepsilon_n$ систему из конечного числа функций, к которым разрешено
добавить все постоянные $a\in\R$.

\textbf{Проба.} Система $H_\varepsilon$ (1.9) полна по отношению
к $\varepsilon_n$: функции из $\varepsilon_n$ элементарны, т.\,е. по
определению являются суперпозициями функций из $H_\varepsilon$. Но, кроме
констант, $H_\varepsilon$ содержит бесконечно много функций: $a^x$ и $\log_ax$
"--- при каждом $a>0$, $x^n$ "--- при каждом показателе $n$. Константы же нам
даны, и это подсказывает, как сократить систему: $a^x=e^{x\ln a}$, где
$\ln a$ "--- константа, т.\,е. все показательные функции получаются из одной
функции $e^x$ умножением аргумента на константу; точно так же
$\log_ax=\frac1{\ln a}\ln x$.

\textbf{Сведение к (1.9).} Достаточно выразить через $H$ каждую функцию
системы (1.9). Тогда $H_\varepsilon\subset\M(H)$, а так как суперпозиция
суперпозиций есть суперпозиция (решение упражнения~4 \S\,4, шаг~2),
получаем $\M(H_\varepsilon)\subset\M(H)$ и, значит,
$\varepsilon_n\subset\M(H_\varepsilon)\subset\M(H)$. Это рассуждение годится
при любом прочтении $\varepsilon_n$ (см. решение упражнения~3): нужно лишь,
чтобы функции из $\varepsilon_n$ были элементарными.

Выражать надо точно: две функции равны, только если у них одна и та же
область определения. Поэтому договоримся об областях определения функций
(1.9). $x/y$ определена при $y\ne0$. $x^n$ при целом $n\ge0$ определена при
всех $x$, при целом $n<0$ "--- при $x\ne0$, а при нецелом $n$ "--- при $x\ge0$,
если $n>0$, и при $x>0$, если $n<0$ (степени с нецелыми показателями в (1.9)
входят: в \S\,4 элементарной названа и функция $\sqrt{t^2}$). $a^x$ определена
при всех $x$; $\log_ax$ ($a\ne1$) "--- при $x>0$; $\tg x$ "--- при $\cos x\ne0$,
$\ctg x$ "--- при $\sin x\ne0$; $\arcsin x$ и $\arccos x$ "--- при $|x|\le1$;
$\arctg x$ и $\operatorname{arcctg}x$ "--- при всех $x$.

\textbf{Система.} Возьмем
\[
H=\{x+y,\ \ xy,\ \ x^y,\ \ \ln x,\ \ \sin x,\ \ \arcsin x\}\cup\{a:\ a\in\R\},
\]
где $x^y$ "--- степень как функция двух аргументов: она определена при $x>0$
(и любом $y$) и при $x=0$, $y>0$ (тогда $x^y=0$). Значение при $x=0$ нужно,
например, чтобы функция $\sqrt{t^2}=|t|$ была определена и при $t=0$. Через
одну эту функцию выражаются все степени с нецелыми показателями и все
показательные функции.

Проверим функции (1.9) по очереди, каждый раз сверяя область определения.
\begin{itemize}[leftmargin=1.2em,itemsep=2pt]
\item $x-y=x+(-1)\cdot y$ "--- при всех $x$, $y$.
\item $1/y=y\cdot(y^2)^{-1}$: при $y\ne0$ число $y^2>0$, степень
  $(y^2)^{-1}=1/y^2$ определена, и произведение равно $1/y$; при $y=0$ степень
  $0^{-1}$ не определена. Область "--- в точности $y\ne0$. Отсюда
  $x/y=x\cdot y\cdot(y^2)^{-1}$.
\item $x^n$: при целом $n\ge1$ "--- произведение $x\cdot x\cdots x$; $x^0=1$
  "--- константа; при целом $n=-m<0$ "--- $(1/x)^m$, область $x\ne0$; при
  нецелом $n$ "--- значение функции $x^y$ при $y=n$, с той же областью, что
  договорено выше.
\item $a^x$ ($a>0$) "--- функция $x^y$ с постоянным основанием $a$ вместо $x$
  и переменной $x$ вместо $y$; определена при всех $x$. В частности, так
  получается $e^x$.
\item $\log_ax=\frac1{\ln a}\cdot\ln x$ ($a>0$, $a\ne1$) "--- при $x>0$.
\item $\cos x=\sin\bigl(x+\frac\pi2\bigr)$ "--- при всех $x$.
\item $\tg x=\sin x\cdot\cos x\cdot(\cos^2x)^{-1}$ "--- при $\cos x\ne0$ (как
  $1/y$ выше); $\ctg x=\cos x\cdot\sin x\cdot(\sin^2x)^{-1}$ "--- при $\sin x\ne0$.
\item $\arccos x=\frac\pi2+(-1)\cdot\arcsin x$ "--- при $|x|\le1$.
\item $\arctg x=\arcsin\bigl(x\cdot(1+x^2)^{-1/2}\bigr)$ "--- при всех $x$. Число
  $u=x/\sqrt{1+x^2}$ лежит в $(-1,1)$, так что $\arcsin u\in(-\frac\pi2,\frac\pi2)$,
  и $\tg(\arcsin u)=u/\sqrt{1-u^2}=x$, потому что $1-u^2=1/(1+x^2)$. Далее
  $\operatorname{arcctg}x=\frac\pi2+(-1)\cdot\arctg x$.
\item $x+y$, $xy$, $\sin x$, $\arcsin x$ и константы входят в $H$.
\end{itemize}
(Все эти равенства вместе с областями определения проверены численно
программой \texttt{artifacts/checks/section5.py}, в тысячах точек, включая
точки вне областей определения.)

Итак, $H_\varepsilon\subset\M(H)$, а значит, $\varepsilon_n\subset\M(H)$:
система $H$ полна по отношению к $\varepsilon_n$ и содержит, кроме констант,
шесть функций.

Два замечания о соглашениях. Если степени с рациональными показателями
и нечетными знаменателями считать определенными и при $x<0$ (например,
$\sqrt[3]{x}$ на всей прямой), к $H$ надо добавить еще одну функцию
$\operatorname{sgn}x\cdot|x|^y$ (при $y>0$ она определена при всех $x$): тогда
$x^{p/q}=\bigl(\operatorname{sgn}x\cdot|x|^{1/q}\bigr)^p$. И второе: если
потребовать, чтобы функции самой системы лежали в $\varepsilon_n$ (были
определены всюду), задача меняется по существу "--- исключаются, например,
$\ln x$ и $x^y$; этого требования в определении полноты нет, и мы его
не рассматриваем.

\answer Например, $H=\{x+y,\ xy,\ x^y,\ \ln x,\ \sin x,\ \arcsin x\}$ вместе
со всеми константами: через эти шесть функций выражаются (с теми же
областями определения) все функции системы (1.9), а значит, и все
элементарные функции, в частности все функции из $\varepsilon_n$.
\end{solution}

\begin{solution}[5 (\S\,5)]
\textbf{Как понимать вопрос.} Структура "--- это функциональная схема, как
на рис.~11 и 12: блоки $\varphi_1,\dots,\varphi_s$ (у блока $\varphi_i$
$n_i$ входов), каждый вход соединен с одной из переменных $x_1,\dots,x_n$ или
с выходом другого блока, циклов нет, выход одного из блоков "--- результат.
Если вставить в блоки непрерывные функции $\varphi_i\in C(\R^{n_i})$, схема
вычислит непрерывную функцию $f\in C(\R^n)$ (упражнение~1 при $m=0$).
Отображение $\chi^*$ из условия сопоставляет набору $(\varphi_1,\dots,\varphi_s)$
эту функцию, и спрашивается, совпадает ли его образ с $C(\R^n)$, т.\,е.
можно ли, подбирая функции в блоках, получить \emph{любую} непрерывную функцию
$n$ переменных. Такие структуры назовем \emph{универсальными}. Если, как
в записи $\varphi(x,\varphi(y,\varphi(z,t)))$, в нескольких блоках должна стоять
одна и та же функция, получится еще меньше функций; отрицательные результаты
ниже верны даже тогда, когда функции во всех блоках выбираются независимо.

\textbf{Пробы.} Схема из одного блока $\varphi_1(x_1,\dots,x_n)$ универсальна:
достаточно взять $\varphi_1=f$. Универсальна и всякая схема, в которой есть
блок, получающий на входы все переменные, и путь от его выхода к результату:
в этот блок ставим $f$, а в каждый блок пути "--- функцию, равную тому из
своих аргументов, по которому пришел путь (например, $\varphi(u,v)=u$).

Интереснее схемы, где каждый блок получает лишь часть переменных. Возьмем
при $n=3$ схему $\varphi_1(x,\varphi_2(y,z))$. Многие функции так получаются:
$x+y+z$ ($\varphi_2=y+z$, $\varphi_1(x,u)=x+u$), $xyz$ ($\varphi_2=yz$,
$\varphi_1=xu$), $x\sin(y^2+z)$ и~т.\,д. Попробуем $f=xy+z$. Пусть
$xy+z=\varphi_1(x,u(y,z))$ при всех $x,y,z$, где $u=\varphi_2$. При $x=0$
и при $x=1$ получаем
\[
z=\varphi_1\bigl(0,u(y,z)\bigr),\qquad y+z=\varphi_1\bigl(1,u(y,z)\bigr).
\]
Значит, по значению $u(y,z)$ однозначно восстанавливаются $z$ и $y+z$, т.\,е.
сама точка $(y,z)$: если $u(y,z)=u(y',z')$, то $z=z'$ и $y+z=y'+z'$, откуда
$(y,z)=(y',z')$. Непрерывная функция двух переменных $u$ оказалась
\emph{взаимно однозначной}: разные точки плоскости она переводит в разные
числа. Так не бывает.

\textbf{Лемма.} \emph{Непрерывная функция $u(a,b)$ на любой окружности плоскости
принимает какое-то значение по крайней мере в двух разных точках.}

Пусть окружность имеет центр $(a_0,b_0)$ и радиус $R$; положим
$g(\theta)=u(a_0+R\cos\theta,\ b_0+R\sin\theta)$. Функция $g$ непрерывна,
и $g(\theta+2\pi)=g(\theta)$. Если $g$ постоянна, утверждение верно. Иначе
пусть $g(\theta_1)$ "--- наименьшее, а $g(\theta_2)$ "--- наибольшее значение,
$\theta_1<\theta_2<\theta_1+2\pi$, и $c$ "--- число строго между ними. Точки
$\theta_1$ и $\theta_2$ делят окружность на две дуги: $\theta\in[\theta_1,\theta_2]$
и $\theta\in[\theta_2,\theta_1+2\pi]$. На концах каждой дуги $g$ принимает
значения $g(\theta_1)<c$ и $g(\theta_2)>c$, поэтому по теореме о промежуточном
значении внутри каждой дуги есть точка, где $g=c$. Внутренности дуг не
пересекаются, так что это две разные точки окружности. $\square$

Итак, $xy+z$ не представляется в виде $\varphi_1(x,\varphi_2(y,z))$
с непрерывными $\varphi_1$, $\varphi_2$, и структура $\varphi_1(x,\varphi_2(y,z))$
не универсальна. Это пример, о котором спрашивается в задаче.

\textbf{Общий признак.} Рассуждение использовало одно свойство схемы:
переменные $y$, $z$ влияют на результат только через выход одного блока.
Скажем, что в схеме есть \emph{узкое место}, если переменные можно разбить
на две непустые группы $x'$ и $x''$, причем в группе $x''$ не меньше двух
переменных и они влияют на результат только через выход одного блока $B$,
который сам зависит только от $x''$: $f=\Phi(x',B(x''))$.
\begin{keystmt}
Схема с узким местом не универсальна: она не дает функцию
$f=x'_1x''_1+x''_2$, где $x'_1$ "--- переменная первой группы,
а $x''_1$, $x''_2$ "--- второй.
\end{keystmt}
Действительно, если $B(p)=B(q)$ для двух наборов значений переменных второй
группы, то $f(x',p)=f(x',q)$ при всех $x'$. Подставив $x'=0$, а затем набор
$x'$ с $x'_1=1$ и остальными нулями, получаем $p_2=q_2$ и $p_1+p_2=q_1+q_2$.
Значит, функция $u(a,b)=B(a,b,0,\dots,0)$ взаимно однозначна на плоскости,
а она непрерывна (как суперпозиция непрерывных функций), что противоречит
лемме. Заметим, что от функции $\Phi$ ничего не требовалось.

Признак сразу дает много неуниверсальных структур. Например, при $n\ge3$ не
универсальна никакая формула из функций одного и двух аргументов (и,
возможно, констант), в которую каждая переменная входит по одному разу.
Назовем вход блока \emph{переменным}, если на него подается подформула,
содержащая хоть одну переменную, и рассмотрим ближайший к выходу блок с
двумя переменными входами. Такой блок есть: если бы у каждого блока был
не более чем один переменный вход, то, спускаясь от выхода по переменным
входам, мы прошли бы цепочку блоков, в которой участвует лишь одна
переменная. Между этим блоком и выходом у блоков лишь по одному переменному
входу, поэтому результат зависит от переменных только через выход нашего
блока. Его переменные входы $A$ и $B$ содержат непересекающиеся группы
переменных (каждая переменная входит в формулу один раз), вместе
составляющие все $n\ge3$ переменных; значит, одна из групп содержит не
меньше двух переменных, скажем группа входа $B$, а группа входа $A$ не
пуста. Это узкое место. Итак, в универсальной структуре из функций одного
и двух аргументов при $n\ge3$ некоторые переменные обязательно
используются несколько раз.

\textbf{Структура $\varphi(x,\varphi(y,\varphi(z,t)))$.} Здесь переменные
$y,z,t$ влияют на результат только через выход среднего блока
$B=\varphi(y,\varphi(z,t))$, а $f=\varphi(x,B)$: это узкое место с группами
$\{x\}$ и $\{y,z,t\}$. Поэтому функция $f(x,y,z,t)=xy+z$ (не зависящая от $t$)
так не представляется "--- ни с одной функцией $\varphi$, ни даже с тремя
разными: $\varphi_1(x,\varphi_2(y,\varphi_3(z,t)))$.

\textbf{Одна структура на кубе.} Универсальная структура должна быть
устроена сложнее: переменные в ней используются многократно. Для функций,
непрерывных на кубе, такую структуру дает теорема Колмогорова~[51]
(усиливающая результат Арнольда~[4]; см. \S\,4, пример~2): всякую функцию $G$,
непрерывную на $[0,1]^n$, можно представить в виде
\[
G(y_1,\dots,y_n)=\sum_{q=1}^{2n+1}\chi_q\Bigl(\sum_{p=1}^{n}\psi_{pq}(y_p)\Bigr),
\]
где $\chi_q$ и $\psi_{pq}$ "--- непрерывные функции одного переменного, причем
$\psi_{pq}$ не зависят от $G$. Сумма двух чисел "--- непрерывная функция двух
аргументов, а функция одного переменного "--- функция двух аргументов, не
зависящая от второго. Поэтому правая часть "--- одна и та же для всех $G$
структура из непрерывных функций двух аргументов; каждая переменная
используется в ней $2n+1$ раз.

\textbf{Одна структура на всем $\R^n$.} В упражнении же речь идет о функциях,
непрерывных на всем $\R^n$. Такая функция может неограниченно расти, и никаким
сдвигом и растяжением переменных ее не свести к функции на кубе. Сведем дело к кубу
в два приема: сначала погасим рост функции множителем, а затем сожмем $\R^n$
в куб.

\emph{Гасящий множитель.} Пусть $f\in C(\R^n)$. Положим
$m(R)=\max\{|f(x)|:\ |x_1|\le R,\dots,|x_n|\le R\}$. Эта функция не убывает
и непрерывна (последнее следует из равномерной непрерывности $f$ на
ограниченных кубах). Положим
\[
v(t)=(1+|t|)\bigl(1+m(|t|)\bigr),\qquad g(x)=\frac{f(x)}{v(x_1)\cdots v(x_n)} .
\]
Функция $v$ непрерывна и $v\ge1$. Оценим $g$. Пусть $R=\max_p|x_p|$
и максимум достигается при $p=p^*$. Точка $x$ лежит в кубе
$\max_p|x_p|\le R$, поэтому $|f(x)|\le m(R)$. В знаменателе все множители
не меньше единицы, а множитель $v(x_{p^*})$ равен $(1+R)(1+m(R))$. Значит,
\[
|g(x)|\le\frac{m(R)}{(1+R)(1+m(R))}<\frac1{1+R},
\]
т.\,е. $g(x)\to0$, когда $\max_p|x_p|\to\infty$.

\emph{Сжатие в куб.} Функция $h(t)=\frac12+\frac1\pi\arctg t$ взаимно
однозначно и непрерывно отображает прямую на интервал $(0,1)$, и обратная
функция $h^{-1}(s)=\tg\bigl(\pi(s-\frac12)\bigr)$ тоже непрерывна. Положим
$G(y)=g\bigl(h^{-1}(y_1),\dots,h^{-1}(y_n)\bigr)$ внутри куба $[0,1]^n$ и $G(y)=0$
на его границе. Функция $G$ непрерывна на всем кубе: внутри "--- как
суперпозиция непрерывных функций; если же точка $y$ изнутри куба стремится
к точке границы, то какая-то координата $y_p$ стремится к $0$ или к $1$,
соответствующее $x_p=h^{-1}(y_p)$ "--- к $-\infty$ или $+\infty$, и по оценке
выше $G(y)\to0$.

\emph{Представление.} Применим к $G$ теорему Колмогорова и подставим
$y_p=h(x_p)$. Получим $g(x)=\sum_q\chi_q\bigl(\sum_p\psi_{pq}(h(x_p))\bigr)$, т.\,е.
\begin{equation*}\tag{2}
f(x_1,\dots,x_n)=v(x_1)\cdots v(x_n)\sum_{q=1}^{2n+1}\chi_q\Bigl(\sum_{p=1}^{n}\psi_{pq}\bigl(h(x_p)\bigr)\Bigr),
\end{equation*}
и это равенство верно при всех $x\in\R^n$. Функции $\chi_q$ заданы на отрезках (областях значений внутренних сумм);
продолжим их на всю прямую постоянными, чтобы все блоки были непрерывны на
$\R$. Правая часть (2) "--- одна и та же для всех $f\in C(\R^n)$ структура
из непрерывных функций одного переменного ($\psi_{pq}\circ h$, $\chi_q$, $v$)
и блоков сложения и умножения; от $f$ зависят лишь $\chi_q$ и $v$.
(Оценка $|g(x)|<1/(1+R)$ проверена численно программой
\texttt{artifacts/checks/section5.py}.)

\answer Полного описания универсальных структур мы не получили; доказано
следующее. (1)~Структура с узким местом "--- группой из двух или более (но
не всех) переменных, влияющих на результат только через выход одного
блока, "--- не универсальна. Пример: $\varphi_1(x,\varphi_2(y,z))$ не дает
функцию $xy+z$. По той же причине при $n\ge3$ не универсальна никакая
формула из функций одного и двух аргументов, в которую каждая переменная
входит один раз. (2)~Не всякая непрерывная функция $f(x,y,z,t)$
представляется в виде $\varphi(x,\varphi(y,\varphi(z,t)))$: так не
получается $xy+z$, даже если три функции $\varphi$ разные. (3)~Одной
структуры достаточно: при каждом $n$ все функции из $C(\R^n)$
представляются по формуле (2) через непрерывные функции одного переменного,
сложение и умножение; меняются лишь функции $\chi_q$ и $v$. Это следует из
теоремы Колмогорова для куба после гашения роста и сжатия $\R^n$ в куб.
\end{solution}

\begin{solution}[6 (\S\,5)]
\textbf{Как перенести вопросы.} Пусть $K$ "--- класс функций (например,
$C^m$). Вставляя в блоки структуры функции класса $K$ соответствующего числа
переменных, получаем функции $n$ переменных, и вопросы упражнения~5
ставятся так: для каких структур так получаются все функции класса
$K(\R^n)$; какие структуры этого не позволяют; существует ли одна структура,
достаточная для представления функций трех и более переменных через функции
двух переменных того же класса. Для классов $C^m$ и $C^A$ суперпозиции снова
лежат в том же классе (упражнения~1 и~2), так что $\chi^*$ отображает
$K(\R^{n_1})\times\dots\times K(\R^{n_s})$ в $K(\R^n)$, и вопрос в точности
такой же, как в упражнении~5: совпадает ли образ с $K(\R^n)$. Обозначение
$H^s(\R^n)$ в книге до этого места не определено; ниже мы скажем и о классах
Гёльдера, и о классах Соболева.

\textbf{Отрицательная часть переносится.} В решении упражнения~5 от функций
в блоках требовалась только непрерывность: взаимно однозначной на плоскости
оказывалась непрерывная функция $u$. Функции классов $C^m$ ($0\le m\le\infty$)
и аналитические функции непрерывны, а свидетель $xy+z$ "--- многочлен, лежащий
во всех этих классах. Поэтому структуры с узким местом, в частности
$\varphi_1(x,\varphi_2(y,z))$ и $\varphi(x,\varphi(y,\varphi(z,t)))$, не
универсальны и в классах $C^m(\R^n)$ и $C^A(\R^n)$: функцию $xy+z$ они не
дают ни при каких непрерывных функциях в блоках.

Для классов, которым многочлены не принадлежат (например, функции класса
должны быть ограничены или убывать на бесконечности), подойдет свидетель,
совпадающий с $x'_1x''_1+x''_2$ лишь вблизи нуля: лемма из решения
упражнения~5 верна для любой окружности, поэтому противоречие получается уже
от взаимной однозначности $u$ в малом круге. Пусть $f=\Phi(x',B(x''))$ и
$f=x'_1x''_1+x''_2$ при всех достаточно малых значениях переменных. Если
$B(p)=B(q)$ для малых $p$, $q$, то $f(x',p)=f(x',q)$ при всех $x'$; подставив
$x'=0$ и набор с малым $x'_1=\delta$ и остальными нулями, получаем $p_2=q_2$
и $\delta p_1=\delta q_1$. Значит, $u(a,b)=B(a,b,0,\dots,0)$ взаимно однозначна
в малом круге "--- противоречие. Такой свидетель есть в любом классе,
содержащем бесконечно дифференцируемые функции, равные нулю вне
ограниченного множества: умножим $x'_1x''_1+x''_2$ на такую функцию, равную
единице около нуля. Непрерывность же функций в блоках нужна для применения
леммы. Поэтому вывод верен для классов Гёльдера и для классов Соболева
$H^s$, когда функции в блоках непрерывны (по теореме вложения
$H^s(\R^k)\subset C(\R^k)$ при $s>k/2$).

\textbf{Положительная часть: проба.} Для непрерывных функций одна
структура (2) из решения упражнения~5 дает все функции. Возможно ли что-либо
подобное для гладких функций? Гладкая функция вблизи точки во многом
описывается своими производными в этой точке, и естественно сосчитать,
сколько чисел определяют производные суперпозиции. Возьмем структуру
$f=\varphi_1(x,\varphi_2(y,z))$ с функциями $\varphi_1,\varphi_2\in C^\infty(\R^2)$.
По формуле (1) решения упражнения~1, примененной несколько раз, производные
$f$ в нуле порядка не выше $N$ выражаются через производные $\varphi_2$
порядка не выше $N$ в точке $(0,0)$ и производные $\varphi_1$ порядка не
выше $N$ в точке $(0,\varphi_2(0,0))$. Производных порядка не выше $N$
(считая саму функцию производной нулевого порядка) у функции двух
переменных столько же, сколько одночленов $u^iv^j$ с $i+j\le N$, т.\,е.
$\binom{N+2}2=\frac{(N+1)(N+2)}2$, а у функции трех переменных
$\binom{N+3}3=\frac{(N+1)(N+2)(N+3)}6$. Сравним:
\begin{center}
\begin{tabular}{c|ccccccc}
$N$&0&1&2&3&4&5&6\\\hline
производные $\varphi_1$ и $\varphi_2$: $2\binom{N+2}2$&2&6&12&20&30&42&56\\
производные $f$: $\binom{N+3}3$&1&4&10&20&35&56&84
\end{tabular}
\end{center}
Начиная с $N=4$ производных у $f$ больше, чем чисел, через которые они
выражаются, и разрыв растет: $\binom{N+3}3$ растет как $N^3$, а
$2\binom{N+2}2$ "--- как $N^2$. Для любой структуры из $s$ блоков не более
чем с двумя входами чисел-<<параметров>> не больше $s\binom{N+2}2$, а это
меньше $\binom{N+3}3$, как только $N+3>3s$. Похоже, что функций трех
переменных <<слишком много>>, чтобы все они получались из функций двух
переменных, как бы ни была устроена структура.

Подсчет параметров сам по себе ничего не доказывает: существуют
непрерывные отображения отрезка на квадрат (кривая Пеано), так что <<меньше
параметров>> не значит <<меньше значений>>. Решающее обстоятельство в другом:
производные суперпозиции выражаются через производные блоков
\emph{многочленами с целыми коэффициентами}, а значения многочленов от малого
числа переменных связаны алгебраическими уравнениями. Проведем рассуждение
для трех переменных $x=(x_1,x_2,x_3)$; обозначения мультииндексов "--- как
в решении упражнения~2, $\partial^\alpha=\partial^{|\alpha|}/\partial
x_1^{\alpha_1}\partial x_2^{\alpha_2}\partial x_3^{\alpha_3}$.

\textbf{Лемма 1 (производные суперпозиции).} \emph{Пусть $f$ получена по
некоторой структуре из функций $\varphi_1,\dots,\varphi_s$ класса $C^\infty$,
и $P_i$ "--- точка, в которой вычисляется блок $\varphi_i$, когда все
переменные равны нулю. Тогда каждое число $\partial^\alpha f(0)$ "--- многочлен
с целыми коэффициентами от чисел $\partial^\beta\varphi_i(P_i)$ ($i=1,\dots,s$,
$|\beta|\le|\alpha|$).}

Сначала одна композиция. Пусть $w(x)=\varphi(u_1(x),\dots,u_p(x))$, где
$\varphi$ и $u_j$ "--- функции класса $C^\infty$. Докажем индукцией по $k=|\alpha|$,
что $\partial^\alpha w(x)$ есть многочлен с целыми коэффициентами (не
зависящий от $\varphi$, $u_j$ и $x$) от чисел $(\partial^\beta\varphi)(u(x))$
и $\partial^\gamma u_j(x)$ с $|\beta|,|\gamma|\le k$. При $k=0$ это сама
функция $w(x)=\varphi(u(x))$. Шаг: продифференцируем такой многочлен по $x_l$.
По формуле (1) решения упражнения~1
\[
\frac{\partial}{\partial x_l}\bigl[(\partial^\beta\varphi)(u(x))\bigr]
=\sum_{j=1}^p(\partial^{\beta+e_j}\varphi)(u(x))\,\frac{\partial u_j}{\partial x_l}(x),
\]
где $e_j$ "--- мультииндекс с единицей на $j$-м месте и нулями на остальных,
а производная от $\partial^\gamma u_j$ есть $\partial^{\gamma+e_l}u_j$. По
правилу дифференцирования произведения производная многочлена от этих чисел
"--- снова многочлен с целыми коэффициентами от таких же чисел, порядки
которых не больше $k+1$.

Теперь вся структура. Занумеруем блоки в порядке вычисления: входы каждого
блока "--- переменные или выходы блоков с меньшими номерами. Индукцией по
номеру блока докажем, что для выхода $w_i$ блока $\varphi_i$ каждое число
$\partial^\alpha w_i(0)$ "--- многочлен с целыми коэффициентами от чисел
$\partial^\beta\varphi_j(P_j)$, $j\le i$, $|\beta|\le|\alpha|$. Индукция полная:
входы блока могут идти от любых предыдущих блоков. Для входов-переменных
числа $\partial^\alpha x_l(0)$ равны 0 или 1. Если для блоков с номерами меньше
$i$ утверждение доказано, применим к $w_i=\varphi_i(u_1,\dots,u_p)$ (где $u_j$
"--- переменные или выходы предыдущих блоков) доказанное для одной
композиции при $x=0$: $\partial^\alpha w_i(0)$ "--- многочлен с целыми
коэффициентами от чисел $\partial^\beta\varphi_i(P_i)$ и $\partial^\gamma u_j(0)$
с $|\beta|,|\gamma|\le|\alpha|$, а последние по предположению индукции "---
многочлены с целыми коэффициентами от чисел $\partial^\beta\varphi_j(P_j)$.
Подстановка многочленов в многочлен дает многочлен. Выход последнего блока
"--- функция $f$. $\square$

\textbf{Лемма 2 (значения многочленов от малого числа переменных).} \emph{Пусть
$D<E$ и числа $c_1,\dots,c_E$ "--- значения в некоторой точке $u^*\in\R^D$
многочленов $P_1,\dots,P_E$ с рациональными коэффициентами от $D$ переменных.
Тогда найдется ненулевой многочлен $R$ с рациональными коэффициентами от $E$
переменных, для которого $R(c_1,\dots,c_E)=0$.}

Пример: числа $c_1=u^2$, $c_2=u^3$ ($D=1$, $E=2$) при любом $u$ связаны
уравнением $c_1^3-c_2^2=0$. В общем случае уравнение найдем подсчетом. Пусть
степени всех $P_j$ не больше $K$. Рассмотрим одночлены $t^\gamma=t_1^{\gamma_1}\cdots
t_E^{\gamma_E}$ степени $|\gamma|\le d$; их $\binom{d+E}E\ge d^E/E!$. Подставив
$t_j=P_j(u)$, получим многочлены от $u$ степени не выше $Kd$ с рациональными
коэффициентами; они лежат в линейном пространстве над $\mathbb Q$ размерности
$\binom{Kd+D}D\le(Kd+D)^D/D!$ (столько одночленов от $D$ переменных степени не
выше $Kd$). Так как $E>D$, при больших $d$ число $d^E/E!$ больше
$(Kd+D)^D/D!$, многочленов больше размерности пространства, и они линейно
зависимы: $\sum_\gamma\lambda_\gamma P(u)^\gamma\equiv0$ с рациональными
$\lambda_\gamma$, не все из которых нули. Многочлен
$R(t)=\sum_\gamma\lambda_\gamma t^\gamma$ ненулевой (одночлены $t^\gamma$
разные), и $R(P(u))\equiv0$; в частности, $R(c_1,\dots,c_E)=R(P(u^*))=0$.
$\square$

\textbf{Алгебраически независимые числа.} Назовем числа $c_1,c_2,\dots$
\emph{алгебраически независимыми}, если никакой ненулевой многочлен
с рациональными коэффициентами от конечного числа из них не обращается
в нуль. Такие числа существуют, и их можно выбирать по одному. Пусть числа
$c_1,\dots,c_k$ уже выбраны и алгебраически независимы (при $k=0$ ничего не
выбрано). Многочленов $R(t_1,\dots,t_{k+1})$ с рациональными коэффициентами
счетное множество. Если $R\ne0$, запишем $R=\sum_jR_j(t_1,\dots,t_k)\,t_{k+1}^j$;
хотя бы один многочлен $R_j$ ненулевой, и по независимости
$R_j(c_1,\dots,c_k)\ne0$, так что многочлен $R(c_1,\dots,c_k,t)$ от одной
переменной $t$ ненулевой и имеет конечное число корней. Все такие корни
составляют счетное множество, а отрезок $[-1,1]$ несчетен; возьмем
$c_{k+1}\in[-1,1]$ вне этого множества. Тогда числа $c_1,\dots,c_{k+1}$
алгебраически независимы.

\textbf{Функция, которую нельзя получить.} Занумеруем мультииндексы
$\alpha=(\alpha_1,\alpha_2,\alpha_3)$ и сопоставим им алгебраически
независимые числа $c_\alpha\in[-1,1]$. Положим
\[
f(x)=\sum_\alpha\frac{c_\alpha}{\alpha_1!\,\alpha_2!\,\alpha_3!}\,x_1^{\alpha_1}x_2^{\alpha_2}x_3^{\alpha_3}.
\]
Члены ряда по модулю не больше членов ряда для $e^{|x_1|}e^{|x_2|}e^{|x_3|}$,
так что ряд сходится абсолютно при всех $x$. В окрестности любой точки $x^0$
разложение $f$ по степеням $x-x^0$ получается раскрытием скобок в
$x^\alpha=\bigl(x^0+(x-x^0)\bigr)^\alpha$, что законно по абсолютной сходимости
(как в решении упражнения~2). Значит, $f\in C^A(\R^3)$, и
$\partial^\alpha f(0)=c_\alpha$ (степенной ряд можно дифференцировать
почленно).

Допустим, что $f$ получена по некоторой структуре из $s$ блоков "--- функций
класса $C^\infty$ не более чем двух аргументов. Возьмем $N$ с $N+3>3s$. По
лемме~1 $E=\binom{N+3}3$ чисел $c_\alpha=\partial^\alpha f(0)$, $|\alpha|\le N$,
"--- значения многочленов с целыми коэффициентами от чисел
$\partial^\beta\varphi_i(P_i)$, $|\beta|\le N$. У блока с двумя входами таких
чисел $\binom{N+2}2$, с одним входом "--- $N+1$, у константы "--- одно, так
что всего их $D\le s\binom{N+2}2<E$. По лемме~2 числа $c_\alpha$
удовлетворяют нетривиальному уравнению с рациональными коэффициентами, а это
противоречит их алгебраической независимости. Значит, $f$ не получается
ни по какой структуре.

Итак, существует функция $f\in C^A(\R^3)$, которая не является суперпозицией
функций двух аргументов класса $C^\infty$ (а тем более класса $C^A$, так как
$C^A\subset C^\infty$), какую бы структуру ни взять. В частности, ни одна
структура из функций двух аргументов не универсальна в классах
$C^\infty(\R^3)$ и $C^A(\R^3)$, и положительный результат упражнения~5 на эти
классы не переносится. Тот же подсчет годится для любого $n\ge2$ и блоков
не более чем с $k<n$ входами: $s\binom{N+k}k<\binom{N+n}n$ при больших $N$,
потому что левая часть растет как $N^k$, а правая "--- как $N^n$.
(Таблица и неравенство $s\binom{N+2}2<\binom{N+3}3\iff N+3>3s$ проверены
программой \texttt{artifacts/checks/section5.py}.)

\textbf{Конечная гладкость.} Для $C^m$ с конечным $m$, классов Гёльдера
и Соболева это рассуждение не проходит: оно требует производных порядка
$N>3s-3$, т.\,е. сколь угодно высоких для больших структур, а у функций этих
классов таких производных нет. Суперпозиции функций конечной гладкости
изучал А.\,Г.~Витушкин (см. обзор: А.\,Г.~Витушкин, Г.\,М.~Хенкин, Линейные
суперпозиции функций, Успехи мат. наук, 1967, т.~22); здесь эти классы мы
не разбираем.

\answer (1)~Отрицательные результаты упражнения~5 переносятся на классы
$C^m(\R^n)$ ($0\le m\le\infty$), $C^A(\R^n)$, классы Гёльдера и классы Соболева
$H^s$ с непрерывными функциями: структуры с узким местом, в частности
$\varphi(x,\varphi(y,\varphi(z,t)))$, не универсальны и там (свидетель "---
$xy+z$ или функция, совпадающая с $xy+z$ вблизи нуля). (2)~Положительный
результат не переносится на $C^\infty$ и $C^A$, причем в сильной форме:
существует функция $f\in C^A(\R^3)$, не являющаяся суперпозицией функций
двух аргументов класса $C^\infty$ ни по какой структуре. Поэтому в этих
классах функциями двух аргументов не обойтись ни при какой структуре.
(3)~Классы конечной гладкости мы не разбирали.
\end{solution}
